% ============================================================================
% Chapter 1: Introduction --- The Case for Native Quantized Intelligence
% ============================================================================

\chapter{Introduction: The Case for Native Quantized Intelligence}
\label{ch:introduction}

\begin{quote}
\textit{``We train models at 32 bits and then throw away 31 of them. The tragedy
is not waste --- it is the belief that those 31 bits were ever necessary.''}
\end{quote}

% ------------------------------------------------------------------
\section{The Cost of Intelligence}
\label{sec:cost_of_intelligence}

The computation required to train and serve modern neural networks has become
unsustainable along three axes: energy, memory, and silicon time. We establish
concrete numbers motivating a fundamentally more efficient paradigm. The
trajectory from 64M-parameter models to 1.8T-parameter frontier systems spans
six orders of magnitude, yet cost grows super-linearly due to memory wall
effects and distributed training communication overhead.

Table~\ref{tab:training_cost} summarizes requirements across model scales.
All figures assume dense Transformers trained on web-scale corpora using
approximately $20\times$ the parameter count in tokens, the empirical
optimum of Hoffmann et al.~\cite{hoffmann2022}.

\begin{table}[htbp]
\centering
\caption{Computational cost of training Transformer models at various scales.
Energy estimates use PUE $= 1.1$, GPU efficiency $= 0.35$ FLOPS/watt (H100).}
\label{tab:training_cost}
\rowcolors{2}{tableaccent}{white}
\begin{tabular}{lrrrrr}
\toprule
\textbf{Model} & \textbf{Params} & \textbf{FLOP} & \textbf{GPU-Hrs} &
\textbf{Memory} & \textbf{Energy} \\
  & & \textbf{(Train)} & \textbf{(H100)} & \textbf{(FP32)} & \textbf{(MWh)} \\
\midrule
GPT-2 Small      & 124M  & 1.3e+18 & 40        & 0.50 GB    & 0.02  \\
LLaMA-7B         & 7B    & 1.1e+21 & 830       & 28.0 GB    & 1.3   \\
LLaMA-65B        & 65B   & 1.0e+22 & 7{,}700   & 260 GB     & 11.4  \\
LLaMA-2-70B      & 70B   & 1.2e+22 & 9{,}200   & 280 GB     & 13.6  \\
GPT-4 (est.)     & 1.8T  & 2.2e+25 & 250{,}000 & 7{,}200 GB & 370   \\
\bottomrule
\end{tabular}
\end{table}

Training GPT-4-class models requires an estimated 370~MWh --- the annual
consumption of 34 US households. A 7B model demands 1.3~MWh. At 1.8T
parameters the FP32 weight matrix alone consumes 7.2~TB, requiring thousands
of GPUs. Between 2020 and 2025, frontier training compute increased 4$\times$
per year while GPU FLOPS improved only 2$\times$; the gap is widening.

For autoregressive inference, each token requires reading the entire weight
matrix once. At 40~GB/s (consumer DDR5), a 70B FP32 model achieves only
143 tokens/second; at INT4, 1{,}143 tokens/second. Memory bandwidth --- not
arithmetic throughput --- is the binding constraint, making weight precision
the most important lever for inference speed.

% ------------------------------------------------------------------
\section{Post-Training Quantization and Its Limits}
\label{sec:post_training_quantization}

Post-training quantization (PTQ) compresses weight tensors after training
completes. Three method families dominate.

\subsection{GPTQ: Hessian-Aware Quantization}

GPTQ~\cite{frantar2023} applies second-order quantization using the inverse
Hessian to determine optimal quantization grids, achieving 3--4 BPW with
modest degradation. However, it treats quantization as error-minimization on
a frozen weight tensor --- the optimization landscape is fixed, and quality
loss is irreducible because the solution was found in $\mathbb{R}^d$ and
projected into a discrete subspace.

\subsection{AWQ: Activation-Aware Weight Quantization}

AWQ~\cite{ma2023awq} preserves higher precision for the approximately 1\%
of ``salient'' weights by scaling activation magnitudes. It improves upon
uniform methods but still operates post-hoc: the model never learns to route
information through low-precision weights, and the activation scaling is a
heuristic, not a learned adaptation.

\subsection{GGUF: The Ecosystem Format}

GGUF provides practical quantized inference (Q2\_K through Q8\_0) with mixed
block-level strategies. It is inference-only: models must be trained in
FP32/BF16 then converted, with each step introducing quality loss and format
friction.

\subsection{The Fundamental PTQ Problem}

Let $w^* \in \mathbb{R}^d$ be the optimal FP32 weight vector. PTQ seeks
$w_q \in \mathcal{C}^d$ minimizing distortion:

\begin{equation}
\label{eq:ptq_projection}
    w_q = \argmin_{w \in \mathcal{C}^d} \|w - w^*\|^2
\end{equation}

The gap $\|w^* - w_q\|$ is the \emph{projection loss} --- irreducible,
non-zero, and harmful to generalization. PTQ is inherently one-shot: quality
is bounded by the FP32 model it compresses. If that model overfits, PTQ
preserves the overfitting. The mutual information between training data and
weights is maximized for the continuous representation; projection to a
discrete lattice can only decrease or maintain it, never increase it.

% ------------------------------------------------------------------
\section{The Native Quantization Thesis}
\label{sec:native_quantization_thesis}

\OIL{} (Optimal Information Lattice) proposes training \textbf{in the
quantized format from the first gradient step}. This is a different
optimization algorithm with provably different convergence and
generalization properties.

\subsection{Why Native Training is Different}

Weights are codebook indices from initialization. The forward pass
dequantizes through learned codebooks (updated via EMA); the backward pass
uses the Straight-Through Estimator (STE) for gradient propagation.

\begin{mythm}[Quantization Barrier]{thm:quant_barrier_intro}
Let $\mathcal{C}$ be a codebook with minimum gap
$\delta = \min_{i \neq j} |c_i - c_j|$, learning rate $\eta$, and gradient
$g_t$ at step $t$. When $\|g_t\| < \delta / (2\eta)$, the parameter update
is identically zero.
\end{mythm}

The dead zone $[-\delta/(2\eta),\; +\delta/(2\eta)]$ suppresses sub-threshold
gradient noise while preserving high-sensitivity directions. This has no
analogue in FP32 training. The STE is the discrete analogue of the
reparameterization trick from variational autoencoders; codebook EMA updates
provide the learning signal while STE gradients provide the direction.

\subsection{Projection versus Adaptation}

PTQ optimizes in $\mathbb{R}^d$ then projects; \OIL{} optimizes directly in
$\mathcal{C}^d$. PTQ quality is always $\leq$ FP32; \OIL{} quality can be
$\geq$ FP32 under the Critical Importance Distribution (CID) assumption ---
that weight sensitivity follows a power-law with exponent $p \approx
0.8$--$1.2$. Under CID, native \OIL{} training achieves matched or better
generalization than FP32 at 1.5~BPW (Theorem~\ref{thm:empirical_superiority}
in Chapter~\ref{ch:native_training}) --- a generalization result, not a
compression result.

% ------------------------------------------------------------------
\section{The OIL Format System}
\label{sec:oil_format_system}

\OIL{} defines 33 quantization formats spanning 1.0 to 32.0 BPW, each
backed by learned codebooks updated via Lloyd-Max optimization and EMA
(Chapter~\ref{ch:oil_formats}). Formats are ordered by representational
capacity from BINARY (1.0 BPW) to OIL32 (32.0 BPW, native FP32).

\begin{table}[htbp]
\centering
\caption{The OIL format family. GRP variants provide lossless block-grouped
storage with exact roundtrip reconstruction.}
\label{tab:oil_formats}
\rowcolors{2}{tableaccent}{white}
\begin{tabular}{lrrrr}
\toprule
\textbf{Format} & \textbf{BPW} & \textbf{Bytes/Wt} &
\textbf{Info Wt} & \textbf{Variant} \\
\midrule
BINARY      & 1.00  & 0.125   & 32.0$\times$ & Fixed \\
SPARK\_Q0   & 1.50  & 0.1875  & 21.3$\times$ & Single, 2-mix, 4-mix \\
TERNARY     & 1.58  & 0.1975  & 20.3$\times$ & Fixed \{$-s, 0, +s$\} \\
OIL2        & 2.00  & 0.250   & 16.0$\times$ & GRP available \\
OIL4        & 4.00  & 0.500   & 8.0$\times$  & GRP available \\
OIL8        & 8.00  & 1.000   & 4.0$\times$  & GRP available \\
OIL16       & 16.00 & 2.000   & 2.0$\times$  & GRP available \\
OIL32       & 32.00 & 4.000   & 1.0$\times$  & FP32 native \\
\bottomrule
\end{tabular}
\end{table}

Non-GRP formats are lossy. GRP variants store quantized blocks with residual
correction metadata enabling exact roundtrip reconstruction (MSE $= 0.0$)
while maintaining substantially better compression than FP32. Mixed
precision uses the twimix framework (Chapter~\ref{ch:mixed_precision}):
2-mix and 4-mix CID-weighted allocation. Trimix is explicitly disallowed.

% ------------------------------------------------------------------
\section{SOPS: A New Metric}
\label{sec:sops_metric}

FLOPS is blind to bit-width: it counts multiply-adds identically whether
operands occupy 32 bits or 2 bits, rendering it meaningless for comparing
quantized and full-precision systems. On a 40~GB/s memory subsystem, FP32
weights (4 bytes) yield $10 \times 10^9$ multiply-adds/s while OIL4 weights
(0.5 bytes) yield $80 \times 10^9$ --- FLOPS reports equal counts despite
8$\times$ efficiency difference.

\SOPS{} (Sextillion Operations Per Second) counts information-weighted
operations. The conservation law (Theorem~\ref{thm:conservation} in
Chapter~\ref{ch:sops}) provides:

\begin{mythm}[SOPS Conservation Law]{thm:conservation_preview}
Every byte of weight memory, regardless of format, yields exactly 4
FP32-equivalent operations.
\end{mythm}

This exact identity holds across the entire BPW spectrum, providing a
principled basis for cross-precision comparison.

% ------------------------------------------------------------------
\section{The MYTHOS.cpp Engine}
\label{sec:mythos_engine}

\MYTHOS{}.cpp is a zero-dependency C++20 engine: 88{,}000+ lines across
337+ source files implementing tokenization, OIL-format training with STE
backward passes and codebook EMA updates, and inference with hand-written
SIMD kernels (SSE4.2, AVX2, AVX-512). It compiles to 82 targets across
Windows and Linux with zero errors and zero warnings.

\begin{table}[htbp]
\centering
\caption{\MYTHOS{}.cpp build target breakdown by subsystem.}
\label{tab:build_targets}
\rowcolors{2}{tableaccent}{white}
\begin{tabular}{lr}
\toprule
\textbf{Subsystem} & \textbf{Targets} \\
\midrule
Core library (\texttt{libmythos})                          & 1  \\
Format codec (\texttt{liboil\_codec})                     & 1  \\
Tokenizer engine                                           & 3  \\
Training pipeline                                          & 5  \\
Inference engine                                           & 4  \\
SIMD kernels (SSE4.2, AVX2, AVX-512, Scalar)              & 16 \\
GPU backend (Vulkan compute, SPIR-V)                       & 8  \\
Autograd engine                                            & 4  \\
Distributed training (FSDP, TP, RingAllReduce)             & 6  \\
Memory-mapped loader                                       & 3  \\
Benchmark suite                                            & 12 \\
Test harness                                               & 19 \\
Utilities and tools                                        & 1  \\
\midrule
\textbf{Total}                                             & \textbf{82} \\
\bottomrule
\end{tabular}
\end{table}

No Python, no pip packages, no CUDA toolkit, no shared libraries beyond the
OS. A single statically-linked binary trains and infers on any x86-64
processor with SSE4.2. The Vulkan backend supports AMD, Intel, and mobile
hardware with zero-copy unified memory on integrated GPUs.

% ------------------------------------------------------------------
\section{Summary of Contributions}
\label{sec:contributions}

\begin{enumerate}
    \item \textbf{The OIL Format System} (Chapter~\ref{ch:oil_formats}):
    33 formats spanning 1.0--32.0 BPW with learned codebooks, lossless GRP
    variants, and the quantization barrier theorem.

    \item \textbf{Mixed Precision via Twimix} (Chapter~\ref{ch:mixed_precision}):
    2-mix and 4-mix CID-weighted allocation provably maximizing information
    capacity under memory constraints.

    \item \textbf{SOPS Metric} (Chapter~\ref{ch:sops}): Information-weighted
    operations with an exact conservation law: every byte yields 4
    FP32-equivalent operations.

    \item \textbf{Native Training Without FP32 Master Weights}
    (Chapter~\ref{ch:native_training}): STE-based training with exponential
    convergence, PAC-Bayes bounds, and diagonal dominance proofs.

    \item \textbf{Continual Learning} (Chapter~\ref{ch:continual_learning}):
    Stability-plasticity equilibrium with EWC regularization and forgetting
    benchmarks.

    \item \textbf{Multi-Token Prediction} (Chapter~\ref{ch:mtp}): Parallel
    decoding with speculative verification and KV rollback for 2--4$\times$
    speedup.

    \item \textbf{Lock-Free Thread Safety} (Chapter~\ref{ch:thread_safety}):
    Wait-free data structures enabling linear scaling to 64+ cores.

    \item \textbf{\MYTHOS{}.cpp Engine} (Chapter~\ref{ch:empirical_results}):
    88{,}000+ lines of zero-dependency C++20, 337+ files, 82 targets.

    \item \textbf{GPU Backend} (Chapter~\ref{ch:gpu_backend}): Vulkan compute
    with SPIR-V shaders, iGPU zero-copy, DirectX 12 support.

    \item \textbf{Formal Theory} (Chapters~\ref{ch:oil_formats}--\ref{ch:native_training}):
    13 theorems with complete proofs including the quantization barrier,
    conservation law, and PAC-Bayes bounds.
\end{enumerate}

% ------------------------------------------------------------------
\section{Paper Organization}
\label{sec:organization}

Chapter~\ref{ch:oil_formats} defines all OIL formats, codebooks, GRP
variants, and proves the quantization barrier theorem. Chapter~\ref{ch:mixed_precision}
develops twimix 2-mix and 4-mix allocation with CID theory.
Chapter~\ref{ch:sops} establishes SOPS with the conservation law.
Chapter~\ref{ch:native_training} presents native STE training with
convergence proofs and PAC-Bayes bounds. Chapter~\ref{ch:continual_learning}
addresses stability-plasticity for lifelong adaptation. Chapter~\ref{ch:mtp}
covers multi-token prediction and speculative decoding. Chapter~\ref{ch:thread_safety}
details lock-free data structures. Chapter~\ref{ch:empirical_results}
presents benchmarks across 40 seeds and 4 model scales. Chapter~\ref{ch:gpu_backend}
describes the Vulkan compute backend. Chapter~\ref{ch:distributed} covers
tensor parallelism, FSDP, and RingAllReduce. Chapter~\ref{ch:future_directions}
discusses hardware acceleration and roadmap. Chapter~\ref{ch:conclusion}
summarizes contributions.

Readers focused on theory should read
Chapters~\ref{ch:oil_formats}--\ref{ch:native_training}. Systems engineers
should read Chapters~\ref{ch:thread_safety}--\ref{ch:gpu_backend}.

% ------------------------------------------------------------------
\section{Democratizing AI Research}
\label{sec:democratization}

The AI ecosystem has a structural access problem. Training a competitive model
requires Python with PyTorch and CUDA (4--15~GB installation), NVIDIA GPUs
(CUDA lock-in), and distributed infrastructure (8--1024 GPUs for weeks).

\textbf{Financial.} A single H100 costs \$30{,}000. A 7B training run costs
\$1{,}000--\$5{,}000 on cloud providers. Frontier models cost millions,
exceeding the annual budget of most universities.

\textbf{Infrastructure.} A ``simple'' inference server requires a Python
environment with 847 transitive dependencies, a 6~GB Docker image, and a
running process. Air-gapped hospitals, embedded devices, and financial
institutions requiring deterministic inference cannot accommodate this.

\textbf{Supply chain.} A \texttt{pip install torch} pulls in 300+ unvetted
transitive dependencies --- each an attack vector for weight exfiltration or
backdoor injection.

\MYTHOS{}.cpp eliminates all three barriers. A single statically-linked C++20
binary performs training and inference with zero dependencies. A 14-year-old
developer on a Ryzen 5 5600GT with 14~GB RAM built the entire system. The
barrier to production AI infrastructure is cultural, not computational.

\begin{table}[htbp]
\centering
\caption{Framework comparison. \OIL{} is the only system supporting sub-byte
native training, lossless GRP formats, and zero-dependency deployment.}
\label{tab:framework_comparison}
\rowcolors{2}{tableaccent}{white}
\begin{tabular}{lcccc}
\toprule
\textbf{Feature} & \textbf{\MYTHOS{}.cpp} & \textbf{llama.cpp} &
\textbf{BitNet.cpp} & \textbf{MLX} \\
\midrule
Language         & Pure C++20    & C/C++        & C++         & Python+Metal \\
Dependencies     & None          & ggml         & PyTorch     & Swift        \\
Native Training  & Yes (\OIL{})  & No           & Yes (tern.) & Yes (FP16)   \\
Min BPW          & 1.0           & 1.5 (Q2\_K)  & 1.58        & 16 (FP16)    \\
Mixed Precision  & Yes (twimix)  & Yes (QAT)    & No          & No           \\
Lossless Formats & Yes (GRP)     & No           & No          & No           \\
SOPS Metric      & Yes           & No           & No          & No           \\
Sub-byte Train   & Yes           & No           & No          & No           \\
Build Targets    & 82            & $\sim$30     & $\sim$5     & $\sim$10     \\
iGPU Zero-Copy   & Yes (Vulkan)  & No           & No          & Yes (Metal)  \\
Distributed      & FSDP, TP      & No           & No          & No           \\
\bottomrule
\end{tabular}
\end{table}

A researcher in a resource-constrained environment can download a single
binary, run \texttt{mythos\_train --format oil4 --epochs 100
--data corpus.bin}, and begin training immediately. The format, optimizer,
SIMD kernels, and inference engine are contained in one zero-dependency
executable. This is the operational reality of the system described in
this paper.

The practical implications extend beyond individual research productivity. By removing the dependency on cloud GPU instances and proprietary software stacks, \MYTHOS{}.cpp enables AI research in contexts where connectivity, budget, or institutional policies prevent access to conventional ML infrastructure. Educational institutions in developing economies, independent researchers without institutional affiliation, and hobbyists exploring neural network internals all benefit from a toolchain that requires nothing more than a C++20 compiler and a commodity computer. This democratization of access to advanced quantization research is not incidental to the project---it is a core design goal, and every architectural decision from the choice of C++20 to the zero-dependency policy serves this objective.
